\section{The Scale of Chords}
Barry's \emph{scale of chords} concept is all about breaking out of the mold of the standard chord shapes and enabling harmonic movement. Instead of seeing each chord as a set of four notes, we can reimagine them as two superimposed chords that create an eight note scale. To create this eight note scale we take our original four notes and add the four notes of a diminished chord starting a step above the root of the chord. This leads to the creation of four scales of chords. We have the major sixth diminished scale, the minor sixth diminished scale, the dominant seventh diminished scale, and the dominant seventh flat five diminished scale.

\bookExample{Major Sixth Diminished Scale of Chords}{./excerpts/major-6th-dim.ly}
\bookExample{Minor Sixth Diminished Scale of Chords}{./excerpts/minor-6th-dim.ly}
\bookExample{Dominant Seventh Diminished Scale of Chords}{./excerpts/dom-7-dim.ly}
\bookExample{Dominant Seventh Flat Five Diminished Scale of Chords}{./excerpts/dom-7-flat-5-dim.ly}

\section{Borrowing}
The concept of \emph{borrowing} is quite similar to the appoggiatura. Looking at the above scale of chords example, we call the diminshed chords, \emph{off chords}. When we talk about borrowing notes, we are talking about playing a single chord whose notes belong to the on chord \emph{and} the off chord. For example, we can play the \MusicChord{root=C, quality=maj6} chord and borrow the \mnote{B} from the \MusicChord{root=D, quality=dim}. We can even do this in sequence like so:

\bookExample{Borrowing a Notes from the Off Chord}{./excerpts/borrowing-basics.ly}

Instead of borrowing a single note we can also borrow two notes from the other chord:

\bookExample{Borrowing Two Notes from the Off Chord}{./excerpts/borrowing-double.ly}

Here is a beautiful Barry Harris exercise based on borrowing:

\bookExample{Sequetial Borrowing}{./excerpts/seq-borrowing.ly}

\section{Musical Genesis}
Barry would often go into a whimsical explanation of where music comes from. He would say that in the beginning God created the twelve notes of the chromatic scale. And, from that scale He divided it in two to create the two whole tone scales. The two whole tone scales are \emph{man} and \emph{woman}. We can see in each diminished chord that, two of the notes come from one whole tone scale, and the other two come from the other whole tone scale. When you take the diminished chord and lower a note, you get a dominant chord. If you do this for each of the four notes, you will see that a diminished chord contains four dominant chords. These chords are said to be related, they're \emph{family}. If you raise a note, you get a minor seventh flat five. The minor seventh flat five chords are the what barry called the \emph{minor sixth on the fifth} of each of the dominant chords.

\bookExample{Chromatic Scale}{./excerpts/chromatic-scale.ly}

\bookExample{The Two Whole Tone Scales}{./excerpts/whole-tone.ly}

\bookExample{The Three Diminished Chords}{./excerpts/three-dims.ly}

\bookExample{The Four Dominants From Each Diminished Chord} {./excerpts/four-doms.ly}

\bookExample{The Four Minor Sixths From Each Diminished Chord} {./excerpts/four-min6s.ly}
